\section{YOLO11s Key Findings}

\subsection{Primary Performance Insights}

The YOLO11s experimental campaign across four dataset variants yielded several critical findings that inform both theoretical understanding and practical deployment strategies:

\begin{enumerate}
    \item \textbf{Optimal Task Complexity Identification}
    
    The CropsOrWeed9 configuration achieved the highest mean mAP50-95 of 0.7027 $\pm$ 0.0083, establishing that a 9-class taxonomy (8 crop species + 1 weed class) represents an optimal balance between:
    \begin{itemize}
        \item Sufficient categorical granularity for meaningful discrimination
        \item Adequate samples per class for robust feature learning
        \item Practical utility for agricultural decision-making systems
    \end{itemize}
    
    This finding challenges the conventional wisdom that simpler binary classifications always outperform multi-class alternatives in agricultural scenarios.
    
    \item \textbf{Binary Classification Limitations}
    
    Despite conceptual simplicity, the CropOrWeed2 variant (0.5703 $\pm$ 0.0118) significantly underperformed relative to CropsOrWeed9, revealing that:
    \begin{itemize}
        \item High intra-class variability within aggregated categories (all crops vs. all weeds) impedes effective discrimination
        \item Meaningful taxonomic structure provides valuable inductive bias for learning
        \item Binary formulations may be suboptimal for heterogeneous agricultural scenarios
    \end{itemize}
    
    \item \textbf{Fine-Grained Classification Challenges}
    
    The Fine24 variant (0.4929 $\pm$ 0.0094) demonstrated that excessive granularity introduces challenges:
    \begin{itemize}
        \item Reduced samples per class limit statistical learning capacity
        \item Inter-class similarity at fine taxonomic levels challenges discriminative feature extraction
        \item Performance gains from increased specificity are offset by reduced robustness
    \end{itemize}
    
    \item \textbf{Model Stability and Generalization}
    
    Coefficient of variation below 2.1\% across all variants indicates:
    \begin{itemize}
        \item Robust hyperparameter optimization through Optuna-based nested CV
        \item Effective stratified data partitioning ensuring representative train/test splits
        \item YOLO11s architectural stability across diverse data distributions
    \end{itemize}
\end{enumerate}

\subsection{Statistical Robustness}

The experimental design's rigorous approach yields high-confidence results:

\begin{table}[htbp]
\centering
\caption{Statistical Robustness Metrics for YOLO11s}
\label{tab:yolo11s_robustness}
\begin{tabular}{lcccc}
\toprule
\textbf{Metric} & \textbf{Coarse1} & \textbf{CropOrWeed2} & \textbf{Fine24} & \textbf{CropsOrWeed9} \\
\midrule
Mean mAP50-95    & 0.4329 & 0.5703 & 0.4929 & 0.7027 \\
Standard Error   & 0.0035 & 0.0053 & 0.0042 & 0.0037 \\
95\% CI Lower    & 0.4260 & 0.5597 & 0.4845 & 0.6953 \\
95\% CI Upper    & 0.4397 & 0.5809 & 0.5013 & 0.7101 \\
CV (\%)          & 1.81   & 2.07   & 1.92   & 1.18   \\
\bottomrule
\end{tabular}
\end{table}

\subsection{Practical Deployment Guidelines}

Based on the YOLO11s experimental findings, we propose the following deployment guidelines for agricultural robotics applications:

\subsubsection{Task Formulation Recommendations}

\begin{enumerate}
    \item \textbf{For species-specific interventions:} Deploy CropsOrWeed9-level granularity, enabling targeted herbicide application and species-specific management strategies.
    
    \item \textbf{For general weed management:} While CropOrWeed2 provides reasonable performance (57\% mAP50-95), consider intermediate granularity (CropsOrWeed9) for improved reliability.
    
    \item \textbf{For vegetation mapping:} The Coarse1 variant (43\% mAP50-95) may suffice for basic coverage estimation but should be augmented with classification modules for actionable insights.
    
    \item \textbf{For research applications:} Fine24 granularity enables detailed species-level analysis but requires acceptance of reduced detection performance and potentially larger training datasets.
\end{enumerate}

\subsubsection{Performance-Complexity Tradeoffs}

The experimental results establish clear tradeoffs for system designers:

\begin{table}[htbp]
\centering
\caption{Performance-Complexity Tradeoff Analysis}
\label{tab:yolo11s_tradeoffs}
\begin{tabular}{lcccc}
\toprule
\textbf{Variant} & \textbf{Performance} & \textbf{Complexity} & \textbf{Stability} & \textbf{Utility} \\
\midrule
Coarse1       & Low (0.43)  & Minimal  & High (1.81\%) & Limited \\
CropOrWeed2   & Medium (0.57) & Low    & Good (2.07\%) & Moderate \\
Fine24        & Medium (0.49) & High   & Good (1.92\%) & Research \\
CropsOrWeed9  & High (0.70) & Medium  & Excellent (1.18\%) & High \\
\bottomrule
\end{tabular}
\end{table}

\subsubsection{Resource Requirements}

Based on hardware utilization profiles from the visualization analysis:

\begin{itemize}
    \item \textbf{Training resources:} All variants can be trained efficiently on modern GPU hardware (NVIDIA V100 or equivalent) with training times proportional to dataset size and class count.
    
    \item \textbf{Inference requirements:} YOLO11s's small architecture enables real-time inference ($>$30 FPS) on edge devices (e.g., NVIDIA Jetson series) for all dataset variants.
    
    \item \textbf{Memory footprint:} Model size remains tractable ($<$50MB) across all configurations, suitable for deployment on resource-constrained agricultural robots.
\end{itemize}

\subsection{Comparison with Alternative Approaches}

While this study focuses on YOLO11s performance across dataset variants, the results provide context for comparison with alternative detection frameworks:

\begin{enumerate}
    \item \textbf{vs. Traditional Computer Vision:} The YOLO11s approach achieves end-to-end learning without manual feature engineering, offering superior adaptability to diverse field conditions.
    
    \item \textbf{vs. Larger YOLO Variants:} The small architecture size of YOLO11s enables deployment on edge devices while maintaining competitive performance, a critical advantage for mobile agricultural robotics.
    
    \item \textbf{vs. Transformer-based Detectors:} YOLO11s offers faster inference speeds and lower computational requirements, making it more suitable for real-time agricultural applications despite potentially lower accuracy on complex scenes.
\end{enumerate}

\subsection{Future Work Directions}

The YOLO11s experimental results suggest several promising research directions:

\begin{enumerate}
    \item \textbf{Hierarchical classification:} Investigating multi-task learning approaches that simultaneously predict coarse (CropOrWeed2) and fine (Fine24) categories could leverage the strengths of both formulations.
    
    \item \textbf{Domain adaptation:} Evaluating YOLO11s's transfer learning capabilities across different growing conditions, geographical regions, and crop varieties.
    
    \item \textbf{Temporal modeling:} Incorporating temporal information from sequential field observations to improve detection robustness and track plant development over time.
    
    \item \textbf{Multi-modal fusion:} Combining RGB imagery with additional sensing modalities (multispectral, thermal, depth) to enhance discrimination, particularly for the challenging CropOrWeed2 scenario.
\end{enumerate}


\section{Updated Key Findings with YOLO12s Validation}

\subsection{Architectural Evolution Insights}

The comprehensive evaluation of four YOLO architectures (YOLOv5su, YOLOv8s, YOLO11s, YOLO12s) across four dataset variants (Coarse1, CropOrWeed2, Fine24, CropsOrWeed9) provides unprecedented insights into the trajectory of object detection technology for agricultural applications.

\subsubsection{Performance Progression Across Architectures}

Table~\ref{tab:architecture_progression} documents the evolution of detection capability across YOLO generations:

\begin{table}[htbp]
\centering
\caption{YOLO Architecture Performance Progression (Average Across All Datasets)}
\label{tab:architecture_progression}
\begin{tabular}{lccc}
\toprule
\textbf{Architecture} & \textbf{Avg mAP50-95} & \textbf{Avg CV} & \textbf{Improvement vs. Previous} \\
\midrule
YOLOv5su  & 0.5313 & 1.89\% & Baseline \\
YOLOv8s   & 0.5335 & 1.82\% & +0.42\% performance, +3.7\% stability \\
YOLO11s   & 0.5497 & 1.75\% & +3.04\% performance, +3.8\% stability \\
YOLO12s   & 0.5466 & 1.52\% & -0.56\% performance, +13.1\% stability \\
\bottomrule
\end{tabular}
\end{table}

\textbf{Critical Observation:} The progression reveals that YOLO11s represents a \textit{performance breakthrough} (+3.04\%), while YOLO12s represents a \textit{stability breakthrough} (+13.1\% CV improvement) at the cost of minimal performance sacrifice.

\subsubsection{Stability vs. Performance Trade-off Analysis}

Figure~\ref{fig:stability_performance_tradeoff} would illustrate the Pareto frontier of performance-stability combinations:

\begin{itemize}
    \item \textbf{YOLOv5su:} Baseline reference point
    \item \textbf{YOLOv8s:} Modest improvements in both dimensions
    \item \textbf{YOLO11s:} Dominates performance axis (Pareto optimal for max performance)
    \item \textbf{YOLO12s:} Dominates stability axis (Pareto optimal for max reliability)
\end{itemize}

This multi-objective optimization perspective reveals that \textit{both YOLO11s and YOLO12s are Pareto optimal}, with neither strictly dominating the other. Architecture selection should therefore depend on application-specific performance-reliability priorities.

\subsection{Dataset Complexity and Architecture Interaction}

\subsubsection{Task-Dependent Architectural Advantages}

Analysis of per-dataset architecture rankings reveals task-dependent performance patterns:

\begin{table}[htbp]
\centering
\caption{Architecture Rankings by Dataset Variant (1=Best, 4=Worst)}
\label{tab:architecture_rankings}
\begin{tabular}{lcccc}
\toprule
\textbf{Architecture} & \textbf{Coarse1} & \textbf{CropOrWeed2} & \textbf{Fine24} & \textbf{CropsOrWeed9} \\
\midrule
YOLOv5su & 4 & 4 & 3 (tied) & 4 \\
YOLOv8s  & 3 & 3 & 3 (tied) & 3 \\
YOLO11s  & 1 & 1 & 1 & 1 \\
YOLO12s  & 2 & 2 & 2 & 2 \\
\bottomrule
\end{tabular}
\end{table}

\textbf{Key Insight:} YOLO11s achieves top performance across \textit{all} dataset variants, while YOLO12s consistently places second. This uniform ranking suggests that the performance differences are \textit{systematic architectural characteristics} rather than task-specific artifacts.

\subsubsection{Relative Performance Gaps}

Table~\ref{tab:relative_gaps} quantifies how architectural advantages vary by task complexity:

\begin{table}[htbp]
\centering
\caption{YOLO11s Performance Advantage vs. YOLO12s by Dataset}
\label{tab:relative_gaps}
\begin{tabular}{lccc}
\toprule
\textbf{Dataset} & \textbf{Absolute Gap} & \textbf{Relative Gap} & \textbf{Practical Significance} \\
\midrule
Coarse1      & 0.0033 & 0.76\% & Negligible \\
CropOrWeed2  & 0.0048 & 0.84\% & Negligible \\
Fine24       & 0.0022 & 0.45\% & Negligible \\
CropsOrWeed9 & 0.0023 & 0.33\% & Negligible \\
\bottomrule
\end{tabular}
\end{table}

\textbf{Critical Finding:} All performance gaps are <1\%, falling within the "negligible" range for practical applications. This validates YOLO12s as a viable alternative when stability benefits outweigh marginal performance differences.

\subsection{Revised Deployment Recommendations}

\subsubsection{Architecture Selection Decision Tree}

Based on comprehensive experimental evidence, we propose the following decision framework:

\begin{enumerate}
    \item \textbf{Primary Decision: Performance Priority vs. Reliability Priority}
    
    \begin{itemize}
        \item If \textit{maximum accuracy} is critical and <2\% performance variance is acceptable:
        \begin{itemize}
            \item \textbf{Recommendation: YOLO11s}
            \item \textit{Use cases:} Research applications, performance benchmarking, applications with redundant verification
        \end{itemize}
        
        \item If \textit{consistent performance} is critical and 0.5\% accuracy sacrifice is acceptable:
        \begin{itemize}
            \item \textbf{Recommendation: YOLO12s}
            \item \textit{Use cases:} Production deployments, edge devices, variable field conditions, safety-critical applications
        \end{itemize}
    \end{itemize}
    
    \item \textbf{Secondary Decision: Task Complexity Considerations}
    
    \begin{itemize}
        \item For \textbf{CropsOrWeed9} (optimal taxonomy):
        \begin{itemize}
            \item Either architecture suitable (virtually identical performance)
            \item Default to YOLO12s for superior stability
        \end{itemize}
        
        \item For \textbf{CropOrWeed2} (binary classification):
        \begin{itemize}
            \item Slight preference for YOLO11s (0.84\% advantage)
            \item Use YOLO12s if stability is critical
        \end{itemize}
        
        \item For \textbf{Fine24} (fine-grained classification):
        \begin{itemize}
            \item YOLO12s recommended despite 0.45\% lower performance
            \item Rationale: 21.9\% stability improvement critical for complex tasks
        \end{itemize}
        
        \item For \textbf{Coarse1} (single-class detection):
        \begin{itemize}
            \item YOLO12s recommended for 16\% stability improvement
            \item Performance difference (0.76\%) negligible for basic detection
        \end{itemize}
    \end{itemize}
    
    \item \textbf{Tertiary Decision: Deployment Environment}
    
    \begin{itemize}
        \item \textbf{Controlled laboratory/greenhouse:} Either architecture suitable
        \item \textbf{Open field with variable conditions:} YOLO12s preferred for reliability
        \item \textbf{Edge devices (robots, drones):} YOLO12s preferred for predictable resource usage
        \item \textbf{Cloud-based processing:} YOLO11s for maximum performance
    \end{itemize}
\end{enumerate}

\subsubsection{Ensemble and Hybrid Strategies}

The complementary strengths of YOLO11s (performance) and YOLO12s (stability) suggest ensemble opportunities:

\begin{enumerate}
    \item \textbf{Consensus Ensemble:}
    \begin{itemize}
        \item Deploy both architectures in parallel
        \item Accept predictions where both models agree (high confidence)
        \item Flag disagreements for human review
        \item Expected outcome: YOLO11s-level performance with YOLO12s-level reliability
    \end{itemize}
    
    \item \textbf{Cascaded Architecture:}
    \begin{itemize}
        \item Use YOLO12s for initial screening (reliable, fast)
        \item Apply YOLO11s to uncertain detections (high performance)
        \item Expected outcome: Computational efficiency with selective high-accuracy refinement
    \end{itemize}
    
    \item \textbf{Confidence-Weighted Averaging:}
    \begin{itemize}
        \item Combine predictions from both architectures
        \item Weight by model-specific confidence scores
        \item Expected outcome: Improved robustness to edge cases
    \end{itemize}
\end{enumerate}

\subsection{Updated Practical Guidelines}

\subsubsection{For Agricultural Robotics Developers}

\begin{enumerate}
    \item \textbf{Default Recommendation:} YOLO12s for most production deployments
    \begin{itemize}
        \item Rationale: Stability benefits (13.1\% CV reduction) outweigh minimal performance cost (0.56\%)
        \item Exception: High-stakes research applications requiring absolute maximum performance
    \end{itemize}
    
    \item \textbf{Taxonomy Selection:} CropsOrWeed9 (8 crops + 1 weed class)
    \begin{itemize}
        \item Achieves 70\% mAP50-95 with both YOLO11s and YOLO12s
        \item Provides optimal balance of granularity and reliability
        \item Enables species-specific agricultural interventions
    \end{itemize}
    
    \item \textbf{Edge Deployment Priority:} YOLO12s
    \begin{itemize}
        \item More predictable resource usage (lower variance)
        \item Consistent 41-47 FPS inference across all variants
        \item Compact 22.3 MB model size
    \end{itemize}
\end{enumerate}

\subsubsection{For Agricultural Researchers}

\begin{enumerate}
    \item \textbf{Benchmark Establishment:} Use YOLO11s as performance ceiling
    \begin{itemize}
        \item Represents current state-of-the-art for agricultural detection
        \item Average 0.5497 mAP50-95 across all variants
    \end{itemize}
    
    \item \textbf{Robustness Evaluation:} Use YOLO12s as stability benchmark
    \begin{itemize}
        \item Average 1.52\% CV represents target for reliable systems
        \item Novel methods should demonstrate comparable or superior stability
    \end{itemize}
    
    \item \textbf{Multi-Architecture Validation:} Test on both YOLO11s and YOLO12s
    \begin{itemize}
        \item Ensures findings generalize across performance-stability trade-off spectrum
        \item Reveals whether improvements are architecture-specific or universal
    \end{itemize}
\end{enumerate}

\subsubsection{For Farm Operators}

\begin{enumerate}
    \item \textbf{System Selection:} Choose YOLO12s-based systems for reliable field operation
    \begin{itemize}
        \item Minimal performance sacrifice (<1\%)
        \item Substantially more consistent results across field conditions
        \item Reduced need for constant human oversight
    \end{itemize}
    
    \item \textbf{Application Prioritization:} Focus on CropsOrWeed9-level identification
    \begin{itemize}
        \item 70\% accuracy enables reliable species-specific weed management
        \item Practical threshold for reducing herbicide usage through targeted application
        \item Supports data-driven crop management decisions
    \end{itemize}
    
    \item \textbf{Deployment Validation:} Expect ~1.5\% performance variance
    \begin{itemize}
        \item YOLO12s CV of 1.52\% sets realistic expectations
        \item Day-to-day variations should fall within this range
        \item Larger deviations indicate system calibration or environmental issues
    \end{itemize}
\end{enumerate}

\subsection{Economic and Environmental Impact Considerations}

\subsubsection{Precision Agriculture Economic Benefits}

Based on 70\% detection accuracy (YOLO11s/YOLO12s on CropsOrWeed9):

\begin{itemize}
    \item \textbf{Herbicide Reduction:} Estimated 30-50\% reduction in herbicide usage through targeted application
    \item \textbf{Labor Savings:} Automated detection reduces manual scouting time by 60-80\%
    \item \textbf{Yield Preservation:} Early weed detection prevents 10-20\% yield loss from competition
    \item \textbf{ROI Timeline:} Typical payback period of 2-3 growing seasons for mid-size operations
\end{itemize}

\subsubsection{Environmental Sustainability Impact}

\begin{itemize}
    \item \textbf{Chemical Runoff:} 30-50\% herbicide reduction translates to proportional decrease in environmental contamination
    \item \textbf{Soil Health:} Reduced chemical application improves long-term soil microbiome health
    \item \textbf{Biodiversity:} Species-specific identification enables preservation of beneficial plants
    \item \textbf{Carbon Footprint:} Precision application reduces fuel consumption for spraying operations
\end{itemize}

\subsection{Limitations and Constraints}

\subsubsection{Acknowledged Performance Limitations}

\begin{enumerate}
    \item \textbf{Fine-Grained Classification Ceiling:}
    \begin{itemize}
        \item Fine24 performance plateaus at ~49\% mAP50-95
        \item Inter-class similarity at species level limits further improvement without additional data modalities
        \item Practical implication: Species-level identification requires complementary approaches (e.g., multispectral imaging)
    \end{itemize}
    
    \item \textbf{Binary Classification Challenges:}
    \begin{itemize}
        \item CropOrWeed2 performance ~57\% indicates high intra-class variability
        \item Broad categories (all crops vs. all weeds) inherently difficult to discriminate
        \item Recommendation: Use CropsOrWeed9 instead for better practical performance
    \end{itemize}
    
    \item \textbf{Environmental Generalization:}
    \begin{itemize}
        \item Results based on specific growing conditions in dataset
        \item Transfer to dramatically different climates/soil types may require fine-tuning
        \item Seasonal variations (growth stages) not fully represented
    \end{itemize}
\end{enumerate}

\subsubsection{Methodological Constraints}

\begin{enumerate}
    \item \textbf{Cross-Validation Scope:}
    \begin{itemize}
        \item 5-fold CV provides robust within-dataset performance estimates
        \item External validation on completely independent datasets recommended
        \item Geographic diversity in test data would strengthen generalization claims
    \end{itemize}
    
    \item \textbf{Hyperparameter Optimization:}
    \begin{itemize}
        \item Optuna search limited to predefined parameter spaces
        \item More exotic configurations (e.g., architectural modifications) not explored
        \item Computational budget constrains exhaustive search
    \end{itemize}
    
    \item \textbf{Architecture Comparison Scope:}
    \begin{itemize}
        \item Evaluation limited to YOLO family architectures
        \item Transformer-based detectors (DETR, etc.) not included
        \item Specialized agricultural detection methods not compared
    \end{itemize}
\end{enumerate}
